\documentclass[12pt]{article}

% === Кодировки и язык ===
\usepackage{cmap}
\usepackage[T2A]{fontenc}
\usepackage{mathtext}
\usepackage[utf8]{inputenc}
\usepackage[english, russian]{babel}

% === Цвета ===
\usepackage[dvipsnames]{xcolor}

% === Математика ===
\usepackage{amsmath, amssymb, wasysym, amsthm}
\usepackage{stmaryrd}
\usepackage{proof}

% === Графика и float'ы ===
\usepackage{wrapfig, float}
\usepackage{graphicx}

% === TikZ / PGFPlots ===
\usepackage{pgfplots}
\pgfplotsset{compat=newest}
\usepackage{pgfplotstable}
\usepackage{tikz, tikz-3dplot}
\usetikzlibrary{calc, patterns, angles, quotes, intersections, automata, positioning, arrows.meta}

% === Таблицы ===
\usepackage{tabularx}
\usepackage{xltabular}
\renewcommand\tabularxcolumn[1]{m{#1}}

% === Разметка ===
\usepackage{geometry}
\geometry{margin=0.5in}

% === Списки и прочее ===
\usepackage[shortlabels]{enumitem}
\usepackage{verbatim}

% ==========================================
% === ТЕОРЕМНЫЕ ОКРУЖЕНИЯ (amsthm) ===
% ==========================================
\theoremstyle{plain}
\newtheorem{theorem}{Теорема}[section]
\newtheorem{corollary}{Следствие}[section]
\newtheorem{lemma}{Лемма}[section]
\newtheorem{statement}{Утверждение}[section]
\newtheorem{axiom}{Аксиома}[section]
\newtheorem{law}{Закон}[section]

\theoremstyle{definition}
\newtheorem{definition}{Определение}[section]
\newtheorem{example}{Пример}[section]
\newtheorem{problem}{Задача}
\newtheorem{method}{Метод}[section]
\newtheorem{event}{Событие}[section]

\theoremstyle{remark}
\newtheorem{note}{Заметка}[section]
\newtheorem{property}{Свойство}[section]
\newtheorem{properties}{Свойства}[section]
\newtheorem{principle}{Правило}[section]
\newtheorem{person}{Личность}[section]
\newtheorem*{solution}{Решение}

% ==========================================
% === КРАСИВЫЕ РАМКИ (tcolorbox) ===
% ==========================================
\usepackage{tcolorbox}
\tcbuselibrary{skins, breakable, theorems}

% Базовые стили для разных типов блоков
\tcbset{
    % Стиль для теорем, лемм (Синий - строгий)
    thmstyle/.style={
        enhanced, breakable,
        frame hidden, boxrule=0pt,
        colback=blue!5!white, % Фон: 5% синего, 95% белого
        borderline west={2.5pt}{0pt}{blue!70!black}, % Вертикальная линия слева
        before skip=10pt, after skip=10pt,
        left=4pt, right=4pt, top=4pt, bottom=4pt
    },
    % Стиль для доказательств (фиолетовый - процесс рассуждения)
    proofstyle/.style={
        enhanced, breakable,
        frame hidden, boxrule=0pt,
        colback=Plum!5!white,
        borderline west={2.5pt}{0pt}{Plum!80!black},
        before skip=10pt, after skip=10pt,
        left=4pt, right=4pt, top=4pt, bottom=4pt,
        sharp corners
    },
    % Стиль для определений и аксиом (Зеленый - фундаментальный)
    defstyle/.style={
        enhanced, breakable,
        frame hidden, boxrule=0pt,
        colback=ForestGreen!5!white,
        borderline west={2.5pt}{0pt}{ForestGreen!80!black},
        before skip=10pt, after skip=10pt,
        left=4pt, right=4pt, top=4pt, bottom=4pt
    },
    % Стиль для примеров и задач (Оранжевый - практический)
    exstyle/.style={
        enhanced, breakable,
        frame hidden, boxrule=0pt,
        colback=orange!5!white,
        borderline west={2.5pt}{0pt}{orange!90!black},
        before skip=10pt, after skip=10pt,
        left=4pt, right=4pt, top=4pt, bottom=4pt
    },
    % Стиль для заметок и свойств (Серый - нейтральный)
    remstyle/.style={
        enhanced, breakable,
        frame hidden, boxrule=0pt,
        colback=gray!5!white,
        borderline west={2.5pt}{0pt}{gray!70!black},
        before skip=10pt, after skip=10pt,
        left=4pt, right=4pt, top=4pt, bottom=4pt
    }
}

% Применяем стили к существующим окружениям amsthm
% 1. Строгая математика
\tcolorboxenvironment{theorem}{thmstyle}
\tcolorboxenvironment{lemma}{thmstyle}
\tcolorboxenvironment{corollary}{thmstyle}
\tcolorboxenvironment{statement}{thmstyle}
\tcolorboxenvironment{law}{thmstyle}

% 2. Доказательство
\tcolorboxenvironment{proof}{proofstyle}
\renewcommand{\qed}{}

% 3. Определения
\tcolorboxenvironment{definition}{defstyle}
\tcolorboxenvironment{axiom}{defstyle}

% 4. Практика
\tcolorboxenvironment{example}{exstyle}
\tcolorboxenvironment{problem}{exstyle}
\tcolorboxenvironment{method}{exstyle}
\tcolorboxenvironment{event}{exstyle}

% 5. Заметки и прочее
\tcolorboxenvironment{note}{remstyle}
\tcolorboxenvironment{property}{remstyle}
\tcolorboxenvironment{properties}{remstyle}
\tcolorboxenvironment{principle}{remstyle}
\tcolorboxenvironment{person}{remstyle}
\tcolorboxenvironment{solution}{remstyle}

% === Гиперссылки (почти последним!) ===
\usepackage[
colorlinks=true,
linkcolor=blue,
urlcolor=blue,
citecolor=blue,
bookmarks=true,
bookmarksopen=false
]{hyperref}


\begin{document}
\tableofcontents
\setcounter{tocdepth}{3}
\newpage

%===========================================
%           МЕХАНИКА
%===========================================
\section{Механика}

\begin{definition}
	Механическое движение~--- изменение пространственного положения тела относительно других тел с течением времени.
\end{definition}
\begin{definition}
	При поступательном движении прямая, проведённая через любые две точки внутри тела, остаётся параллельна самой себе.
\end{definition}
\begin{definition}
	При вращательном движении каждая точка тела вращается по своей окружности, центры этих окружностей лежат на одной прямой~--- оси вращения.
\end{definition}
\begin{statement}
	Любое движение~--- суперпозиция поступательного и вращательного движений.
\end{statement}
\begin{definition}
	Колебательное движение~--- движение, повторяющееся с той или иной точностью во времени.
\end{definition}

\subsection{Кинематика}

\begin{definition}
	Кинематика~--- раздел механики, изучающий способы описания движения и связь величин, характеризующих это движение.
\end{definition}

\begin{statement}
	Для описания движения нужны:
	\begin{itemize}
		\item Система отсчёта (тело отсчёта + система координат + часы).
	\end{itemize}
\end{statement}

\subsubsection{Основные понятия}

\begin{definition}
	Траектория~--- кривая, по которой движется тело (точка).
\end{definition}
\begin{definition}
	Путь $s$~--- длина траектории (скалярная величина, $s \geqslant 0$).
\end{definition}
\begin{definition}
	Перемещение $\vec{\Delta r}$~--- вектор из начального положения тела в конечное.
\end{definition}
\begin{definition}
	Средняя скорость: $\vec{V}_{\text{ср}} = \dfrac{\vec{\Delta r}}{\Delta t}$.
	Мгновенная скорость: $\vec{V} = \lim\limits_{\Delta t \to 0} \dfrac{\vec{\Delta r}}{\Delta t} = \dfrac{d\vec{r}}{dt}$.
\end{definition}
\begin{definition}
	Ускорение: $\vec{a} = \lim\limits_{\Delta t \to 0} \dfrac{\Delta \vec{V}}{\Delta t} = \dfrac{d\vec{V}}{dt}$.
\end{definition}

\subsubsection{Равномерное прямолинейное движение (РПД)}

\begin{definition}
	Равномерное прямолинейное движение~--- движение, при котором за любые равные промежутки времени тело проходит одинаковые участки пути, траектория~--- прямая линия. Скорость $\vec{V} = \mathrm{const}$.
\end{definition}

Формулы РПД:
\[
	V = \frac{s}{t}, \qquad s = V \cdot t, \qquad x = x_0 + V_x t.
\]

\subsubsection{Равноускоренное прямолинейное движение (РУД)}

\begin{definition}
	Равноускоренное движение~--- движение с постоянным ускорением: $\vec{a} = \mathrm{const}$.
\end{definition}

Формулы РУД:
\[
	V_x = V_{0x} + a_x t, \qquad
	x = x_0 + V_{0x} t + \frac{a_x t^2}{2}, \qquad
	s = V_{0x} t + \frac{a_x t^2}{2}.
\]

\begin{statement}[Безвременная формула (золотая формула механики)]
	\[
		V^2 = V_0^2 + 2a s \qquad \Longleftrightarrow \qquad s = \frac{V^2 - V_0^2}{2a}.
	\]
\end{statement}
\begin{proof}
	Из $V = V_0 + at$ выражаем $t = \frac{V - V_0}{a}$. Подставляем в $s = V_0 t + \frac{at^2}{2}$:
	\[
		s = V_0 \frac{V - V_0}{a} + \frac{a}{2}\left(\frac{V - V_0}{a}\right)^2
		= \frac{V_0(V - V_0)}{a} + \frac{(V-V_0)^2}{2a}
		= \frac{2V_0(V-V_0) + (V-V_0)^2}{2a}
		= \frac{V^2 - V_0^2}{2a}.
	\]
\end{proof}

\subsubsection{Относительность движения. Преобразование Галилея}

\begin{definition}[Принцип относительности Галилея]
	Во всех инерциальных системах отсчёта механические явления протекают одинаково (законы механики имеют одинаковый вид).
\end{definition}

Пусть система отсчёта $K'$ движется относительно $K$ со скоростью $\vec{V}_0$. Тогда:
\[
	\vec{r} = \vec{r}' + \vec{V}_0 t, \qquad t = t',
\]
\[
	\vec{V} = \vec{V}' + \vec{V}_0 \qquad \text{(закон сложения скоростей Галилея)}.
\]

В общем случае:
\[
	\vec{V}_{\text{абс}} = \vec{V}_{\text{отн}} + \vec{V}_{\text{пер}}.
\]

\subsubsection{Движение тела, брошенного под углом к горизонту}

Тело брошено с высоты $h_0$ под углом $\alpha$ к горизонту со скоростью $V_0$. Ось $x$~--- горизонтальная, ось $y$~--- вертикальная вверх.

\begin{enumerate}
	\item $V_x = V_0 \cos\alpha = \mathrm{const}$, \quad $x = V_0 \cos\alpha \cdot t$.
	\item $V_y = V_0 \sin\alpha - g t$, \quad $y = h_0 + V_0 \sin\alpha \cdot t - \dfrac{g t^2}{2}$.
\end{enumerate}

\paragraph{Уравнение траектории.}
Исключая $t = \dfrac{x}{V_0 \cos\alpha}$:
\[
	y = h_0 + x \tan\alpha - \frac{g x^2}{2 V_0^2 \cos^2\alpha}.
\]
Траектория~--- парабола.

\paragraph{Максимальная высота.} Из $V_y = 0$: $t_{\text{подъёма}} = \dfrac{V_0 \sin\alpha}{g}$.
\[
	H_{\max} = h_0 + \frac{V_0^2 \sin^2\alpha}{2g}.
\]

\paragraph{Время полёта.} Из $y = 0$:
\[
	t_{\text{полёта}} = \frac{V_0 \sin\alpha + \sqrt{V_0^2 \sin^2\alpha + 2g h_0}}{g}.
\]
При $h_0 = 0$: $t_{\text{полёта}} = \dfrac{2 V_0 \sin\alpha}{g}$.

\paragraph{Дальность полёта.} $L = V_0 \cos\alpha \cdot t_{\text{полёта}}$.\\
При $h_0 = 0$: $L = \dfrac{V_0^2 \sin 2\alpha}{g}$. Максимальная дальность при $\alpha = 45^\circ$.

\paragraph{Конечная скорость.}
\[
	V_{\text{к}} = \sqrt{V_0^2 + 2g h_0} \quad \text{(из закона сохранения энергии или прямого вычисления)}.
\]

\subsubsection{Движение по окружности}

\begin{definition}
	Угловая скорость: $\omega = \lim\limits_{\Delta t \to 0} \dfrac{\Delta\varphi}{\Delta t}$, $[\omega] = \dfrac{\text{рад}}{\text{с}}$.
\end{definition}

\begin{statement}
	$V = \omega R$, \qquad $T = \dfrac{2\pi}{\omega} = \dfrac{2\pi R}{V}$, \qquad $\nu = \dfrac{1}{T}$.
\end{statement}

\paragraph{Вывод центростремительного ускорения.}

Рассмотрим тело, движущееся по окружности радиуса $R$ с постоянной скоростью $V$. За малое время $\Delta t$ тело поворачивается на угол $\Delta\varphi$. Вектор скорости тоже поворачивается на тот же угол $\Delta\varphi$, сохраняя модуль $V$. Из равнобедренного треугольника скоростей:
\[
	|\Delta \vec{V}| \approx V \cdot \Delta\varphi \quad (\text{при малом } \Delta\varphi).
\]
Ускорение:
\[
	a = \lim_{\Delta t \to 0} \frac{|\Delta \vec{V}|}{\Delta t} = V \cdot \lim_{\Delta t \to 0} \frac{\Delta\varphi}{\Delta t} = V \omega = \frac{V^2}{R} = \omega^2 R.
\]
Это ускорение направлено к центру окружности (центростремительное).

\paragraph{Угловое ускорение.}
$\beta = \dfrac{d\omega}{dt}$; при $\beta = \mathrm{const}$:
\[
	\omega(t) = \omega_0 + \beta t, \qquad \varphi = \varphi_0 + \omega_0 t + \frac{\beta t^2}{2}.
\]
Тангенциальное ускорение: $a_\tau = \beta R$.\\
Полное ускорение: $a = \sqrt{a_\tau^2 + a_n^2}$, где $a_n = \dfrac{V^2}{R}$.

\subsection{Динамика}

\begin{definition}
	Динамика~--- раздел механики, изучающий причины движения тел (связь движения с действующими силами).
\end{definition}

\subsubsection{Законы Ньютона}

\begin{definition}[Первый закон Ньютона (закон инерции)]
	Существуют такие системы отсчёта (инерциальные), в которых тело, на которое не действуют силы (или действие сил скомпенсировано), сохраняет состояние покоя или равномерного прямолинейного движения.
\end{definition}

\begin{definition}[Второй закон Ньютона]
	\[
		\sum \vec{F} = m \vec{a}.
	\]
	Ускорение тела пропорционально равнодействующей приложенных сил и обратно пропорционально массе тела.
\end{definition}

\begin{definition}[Третий закон Ньютона]
	При взаимодействии двух тел возникают две силы, которые:
	\begin{itemize}
		\item приложены к разным телам,
		\item равны по модулю,
		\item противоположны по направлению,
		\item лежат на одной прямой,
		\item имеют одну природу.
	\end{itemize}
	\[
		\vec{F}_{12} = -\vec{F}_{21}.
	\]
\end{definition}

\begin{statement}[Ограничения применимости]
	Законы Ньютона справедливы:
	\begin{itemize}
		\item при скоростях $V \ll c$ (нерелятивистский случай);
		\item в инерциальных системах отсчёта;
		\item для тел с ненулевой массой покоя.
	\end{itemize}
\end{statement}

\subsubsection{Силы трения}

\begin{definition}
	Сила трения имеет электромагнитную природу. Направлена вдоль поверхности соприкосновения, против относительной скорости движения (или против тенденции к движению).
\end{definition}

\paragraph{Виды трения.}

\begin{enumerate}
	\item \textbf{Трение покоя.} Возникает, когда тело не движется относительно поверхности. Сила трения покоя равна по модулю и противоположна по направлению внешней силе (до предела). Максимальное значение:
	\[
		F_{\text{тр.покоя}}^{\max} = \mu_{\text{п}} N.
	\]

	\item \textbf{Трение скольжения.} Возникает при скольжении тела по поверхности:
	\[
		F_{\text{тр}} = \mu N,
	\]
	где $\mu$~--- коэффициент трения скольжения, $N$~--- сила нормальной реакции опоры.

	\item \textbf{Трение качения.} Значительно меньше трения скольжения. Обусловлено деформацией тела и поверхности:
	\[
		F_{\text{кач}} = \mu_{\text{кач}} \frac{N}{R},
	\]
	где $\mu_{\text{кач}}$~--- коэффициент трения качения (имеет размерность длины), $R$~--- радиус катящегося тела.

	\item \textbf{Вязкое трение} (трение в жидкостях и газах). Пропорционально скорости (при малых скоростях):
	\[
		F_{\text{вязк}} = -\beta V \quad \text{или} \quad F = \frac{\eta V S}{h},
	\]
	где $\eta$~--- динамическая вязкость, $[\eta] = \text{Па}\cdot\text{с}$.\\
	Для вязкого трения \textbf{не существует} силы покоя: при $V = 0$ сила вязкого трения равна нулю.
\end{enumerate}

\subsubsection{Силы упругости. Закон Гука}

\begin{definition}
	Сила упругости~--- сила электромагнитной природы, возникающая при деформации тела и направленная против деформации:
	\[
		F_{\text{упр}} = -k \Delta x \quad \text{(закон Гука)},
	\]
	где $k$~--- жёсткость (коэффициент упругости), $\Delta x$~--- деформация.
\end{definition}

\paragraph{Виды деформаций:}
\begin{itemize}
	\item Упругие (обратимые): растяжение-сжатие, сдвиг, изгиб, кручение.
	\item Пластические (необратимые).
\end{itemize}

\paragraph{Механическое напряжение и модуль Юнга.}
\begin{definition}
	Относительное удлинение: $\varepsilon = \dfrac{\Delta l}{l_0}$.
\end{definition}
\begin{definition}
	Механическое напряжение: $\sigma = \dfrac{F}{S} = E \cdot |\varepsilon|$, $[\sigma] = \text{Па}$.
\end{definition}
\begin{definition}
	Модуль Юнга: $E = \dfrac{k l_0}{S}$, $[E] = \text{Па}$.
\end{definition}

Связь закона Гука с модулем Юнга:
\[
	F = k \Delta l = \frac{E S}{l_0} \Delta l, \qquad \sigma = \frac{F}{S} = E \frac{\Delta l}{l_0} = E |\varepsilon|.
\]

\paragraph{Вывод жёсткости параллельного и последовательного соединения пружин.}

\subparagraph{Параллельное соединение.}
Обе пружины деформируются на одну и ту же величину $\Delta x$. Суммарная сила:
\[
	F = F_1 + F_2 = k_1 \Delta x + k_2 \Delta x = (k_1 + k_2) \Delta x.
\]
Следовательно:
\[
	\boxed{k_{\text{парал}} = k_1 + k_2.}
\]

Через модуль Юнга (две пружины одинаковой длины $l_0$, сечения $S_1$ и $S_2$):
\[
	k = \frac{E S}{l_0} = \frac{E(S_1 + S_2)}{l_0} = k_1 + k_2.
\]

\subparagraph{Последовательное соединение.}
Через обе пружины проходит одна и та же сила $F$. Полная деформация:
\[
	\Delta x = \Delta x_1 + \Delta x_2 = \frac{F}{k_1} + \frac{F}{k_2} = F\left(\frac{1}{k_1} + \frac{1}{k_2}\right).
\]
Поскольку $\Delta x = \dfrac{F}{k_{\text{послед}}}$:
\[
	\boxed{\frac{1}{k_{\text{послед}}} = \frac{1}{k_1} + \frac{1}{k_2}.}
\]

Через модуль Юнга (одинаковое сечение $S$, длины $l_1$ и $l_2$):
\[
	\frac{1}{k} = \frac{l_0}{ES} = \frac{l_1 + l_2}{ES} = \frac{l_1}{ES} + \frac{l_2}{ES} = \frac{1}{k_1} + \frac{1}{k_2}.
\]

\subsubsection{Закон всемирного тяготения}

\begin{law}[всемирного тяготения (Ньютон, 1687)]
	\[
		F = G \frac{m_1 m_2}{R^2}, \qquad G = 6{,}674 \cdot 10^{-11} \; \frac{\text{Н}\cdot\text{м}^2}{\text{кг}^2}.
	\]
\end{law}

\begin{statement}[Границы применимости]
	\begin{itemize}
		\item Точечные тела (или тела, расстояние между которыми много больше их размеров).
		\item Сферически-симметричные тела (плотность зависит только от расстояния до центра).
	\end{itemize}
\end{statement}

\paragraph{Гравитационная и инертная масса.}
Гравитационная масса~--- масса в законе тяготения. Инертная масса~--- масса во втором законе Ньютона. Экспериментально установлено их равенство с высокой точностью (принцип эквивалентности).

\paragraph{Опыт Кавендиша (1798).}
Крутильные весы позволили измерить $G$ (а точнее, плотность Земли). Результат Кавендиша: $\rho_{\text{Земли}} = 5437 \; \text{кг/м}^3$ (современное значение $\approx 5515 \; \text{кг/м}^3$).

\paragraph{Сила тяжести и ускорение свободного падения.}
На поверхности Земли:
\[
	F = mg, \qquad g = \frac{GM_{\text{З}}}{R_{\text{З}}^2} \approx 9{,}81 \; \frac{\text{м}}{\text{с}^2}.
\]
На высоте $h$:
\[
	g(h) = \frac{GM_{\text{З}}}{(R_{\text{З}} + h)^2}.
\]

\paragraph{Вес тела и невесомость.}
\begin{definition}
	Вес тела~--- сила, с которой тело действует на опору или подвес.
\end{definition}

\begin{itemize}
	\item Тело на опоре, ускорение вверх: $P = m(g + a)$.
	\item Тело на опоре, ускорение вниз: $P = m(g - a)$.
	\item При $a = g$ (свободное падение): $P = 0$~--- невесомость.
\end{itemize}

\textbf{Невесомость}~--- состояние, при котором вес тела равен нулю. Физический смысл: тело и опора движутся с одинаковым ускорением (свободное падение), поэтому нет деформации, нет силы реакции. Невесомость наблюдается на орбите, так как спутник~--- это тело в свободном падении (на него действует только гравитация, она обеспечивает центростремительное ускорение).

\paragraph{Центр тяжести.}
Центр тяжести~--- точка, в которой можно считать приложенной силу тяжести. В однородном поле тяжести центр тяжести совпадает с центром масс:
\[
	\vec{r}_c = \frac{\sum_i m_i \vec{r}_i}{\sum_i m_i}.
\]

\paragraph{Первая космическая скорость.}
Тело на круговой орбите вблизи поверхности планеты: гравитационная сила = центростремительная:
\[
	\frac{GMm}{R^2} = \frac{mV_1^2}{R} \quad \Longrightarrow \quad V_1 = \sqrt{\frac{GM}{R}} = \sqrt{gR}.
\]
Для Земли: $V_1 \approx 7{,}9 \; \text{км/с}$.

\paragraph{Законы Кеплера.}

\begin{enumerate}
	\item Все планеты движутся по эллипсам, в одном из фокусов которых находится Солнце.
	\item Радиус-вектор планеты за равные промежутки времени заметает равные площади (закон площадей, следствие сохранения момента импульса).
	\item $\dfrac{T_1^2}{T_2^2} = \dfrac{a_1^3}{a_2^3}$, где $a$~--- большая полуось эллипса.
\end{enumerate}

\begin{proof}[Вывод третьего закона Кеплера для круговых орбит]
	$\frac{GMm}{R^2} = \frac{4\pi^2 m R}{T^2}$, откуда $T^2 = \frac{4\pi^2}{GM} R^3$, т.\,е. $\frac{T^2}{R^3} = \frac{4\pi^2}{GM} = \mathrm{const}$.
\end{proof}

\subsubsection{Импульс тела. Закон сохранения импульса}

\begin{definition}
	Импульс тела: $\vec{p} = m\vec{V}$, $[p] = \text{кг}\cdot\text{м/с}$.
\end{definition}

\begin{definition}[Второй закон Ньютона в импульсной форме]
	\[
		\vec{F} \Delta t = \Delta \vec{p}, \qquad \vec{F} = \frac{d\vec{p}}{dt}.
	\]
\end{definition}

\begin{law}[сохранения импульса]
	Если на систему не действуют внешние силы (или их действие скомпенсировано), то суммарный импульс системы сохраняется:
	\[
		\sum_i \vec{p}_i = \mathrm{const}.
	\]
\end{law}

\begin{proof}
	Рассмотрим систему из двух тел. По третьему закону Ньютона: $\vec{F}_{12} = -\vec{F}_{21}$.
	\[
		\frac{d\vec{p}_1}{dt} = \vec{F}_{21}, \qquad \frac{d\vec{p}_2}{dt} = \vec{F}_{12} = -\vec{F}_{21}.
	\]
	Складывая: $\frac{d}{dt}(\vec{p}_1 + \vec{p}_2) = \vec{0}$, т.\,е. $\vec{p}_1 + \vec{p}_2 = \mathrm{const}$.
\end{proof}

\subsubsection{Реактивное движение}

$\mu$~--- массовый расход топлива ($[\mu] = \text{кг/с}$), $\vec{u}$~--- скорость истечения газов в системе ракеты.

По ЗСИ: $M\vec{V} = (M - \mu\Delta t)(\vec{V} + \Delta\vec{V}) + \mu\Delta t(\vec{V} + \vec{u})$.

Раскрывая и отбрасывая малые второго порядка:
\[
	M\Delta\vec{V} = -\mu\vec{u}\Delta t \quad \Longrightarrow \quad M\vec{a} = -\mu\vec{u}.
\]

\begin{definition}[Уравнение Мещерского]
	$M\vec{a} = \vec{F}_{\text{внеш}} - \mu\vec{u}$, где $\vec{F}_p = -\mu\vec{u}$~--- реактивная сила.
\end{definition}

\begin{definition}[Формула Циолковского]
	Для ракеты без внешних сил, стартующей из покоя:
	\[
		V = u \ln\frac{M_0}{M},
	\]
	где $M_0$~--- начальная масса, $M$~--- конечная масса.
\end{definition}

\subsubsection{Механическая работа. Мощность. Энергия}

\begin{definition}
	Работа постоянной силы: $A = F l \cos\alpha = \vec{F} \cdot \vec{l}$, $[A] = \text{Дж}$.
\end{definition}

\begin{definition}
	Мощность: $P = \dfrac{A}{t} = \vec{F} \cdot \vec{V}$, $[P] = \text{Вт}$.
\end{definition}

\begin{definition}
	Кинетическая энергия: $E_{\text{к}} = \dfrac{mV^2}{2}$.
\end{definition}

\begin{theorem}[О кинетической энергии]
	Работа результирующей силы равна изменению кинетической энергии:
	\[
		A_{\text{рез}} = \Delta E_{\text{к}} = \frac{mV_2^2}{2} - \frac{mV_1^2}{2}.
	\]
\end{theorem}

\begin{definition}
	Потенциальная энергия в поле тяжести: $E_{\text{п}} = mgh$.\\
	Потенциальная энергия пружины: $E_{\text{п}} = \dfrac{k(\Delta x)^2}{2}$.\\
	Потенциальная энергия гравитационного взаимодействия: $E_{\text{п}} = -\dfrac{Gm_1m_2}{R}$.
\end{definition}

\begin{definition}[Консервативные силы]
	Силы, работа которых не зависит от траектории, а определяется только начальным и конечным положениями. Для них: $A = -\Delta E_{\text{п}}$.
\end{definition}

\begin{law}[сохранения механической энергии]
	В замкнутой системе, в которой действуют только консервативные силы:
	\[
		E_{\text{к}} + E_{\text{п}} = \mathrm{const}.
	\]
	Если действуют также неконсервативные силы (трение и т.\,п.):
	\[
		\Delta(E_{\text{к}} + E_{\text{п}}) = A_{\text{неконс}}.
	\]
\end{law}

\begin{definition}[КПД]
	$\eta = \dfrac{A_{\text{пол}}}{A_{\text{зат}}} \cdot 100\%$.
\end{definition}

\subsection{Статика}

\begin{definition}
	Плечо силы~--- кратчайшее расстояние от оси вращения до линии действия силы.
\end{definition}

\begin{definition}
	Момент силы: $M = F \cdot d$, $[M] = \text{Н}\cdot\text{м}$, где $d$~--- плечо.\\
	В векторной форме: $\vec{M} = \vec{r} \times \vec{F}$.
\end{definition}

\begin{definition}[Условия равновесия абсолютно твёрдого тела]
	\begin{enumerate}
		\item Сумма всех сил равна нулю: $\sum \vec{F}_i = \vec{0}$.
		\item Сумма всех моментов сил относительно любой точки равна нулю: $\sum M_i = 0$.
	\end{enumerate}
\end{definition}

\begin{definition}[Центр масс]
	\[
		\vec{r}_c = \frac{\sum_i m_i \vec{r}_i}{\sum_i m_i}.
	\]
\end{definition}

\begin{theorem}[О движении центра масс]
	Центр масс системы движется так, как если бы вся масса была сосредоточена в нём и к нему были приложены все внешние силы:
	\[
		M \vec{a}_c = \vec{F}_{\text{внеш}}.
	\]
\end{theorem}

\paragraph{Виды равновесия.}
\begin{itemize}
	\item \textbf{Устойчивое}~--- при малом отклонении возникает возвращающая сила (центр масс в минимуме потенциальной энергии).
	\item \textbf{Неустойчивое}~--- при малом отклонении равнодействующая уводит от положения равновесия (максимум потенциальной энергии).
	\item \textbf{Безразличное}~--- при малом отклонении равнодействующая остаётся нулевой (потенциальная энергия не меняется).
\end{itemize}

\paragraph{Устойчивость тел.}
Тело на опоре устойчиво, пока вертикаль, проведённая через центр тяжести, проходит через площадь опоры. Чем ниже центр тяжести и шире площадь опоры, тем устойчивее тело.

\subsection{Вращательное движение твёрдых тел}

\subsubsection{Основное уравнение динамики вращательного движения}

\begin{definition}[Момент инерции]
	Для точечного тела: $I = mR^2$.\\
	Для системы точек: $I = \sum_i m_i r_i^2$.\\
	Для сплошного тела: $I = \int r^2 \, dm$.\\
	$[I] = \text{кг}\cdot\text{м}^2$.
\end{definition}

\begin{definition}[Основное уравнение]
	\[
		\sum M = I \beta,
	\]
	где $\beta$~--- угловое ускорение.
\end{definition}

\subsubsection{Моменты инерции тел (с выводом)}

\paragraph{Стержень относительно оси через центр, перпендикулярной стержню.}

Пусть стержень длины $L$, массы $m$, однородный, линейная плотность $\lambda = m/L$. Ось через центр. Элемент $dm = \lambda \, dx$ на расстоянии $x$ от центра:
\[
	I = \int_{-L/2}^{L/2} x^2 \lambda \, dx = \lambda \frac{x^3}{3} \Bigg|_{-L/2}^{L/2} = \lambda \cdot \frac{2(L/2)^3}{3} = \frac{m}{L} \cdot \frac{L^3}{12} = \boxed{\frac{mL^2}{12}}.
\]

\paragraph{Стержень относительно оси через конец.}
По теореме Штейнера: $I = \dfrac{mL^2}{12} + m\left(\dfrac{L}{2}\right)^2 = \dfrac{mL^2}{3}$.

\paragraph{Сплошной однородный шар относительно оси через центр.}
Разобьём шар радиуса $R$ на тонкие диски, перпендикулярные оси. Диск на высоте $z$ имеет радиус $r = \sqrt{R^2 - z^2}$ и толщину $dz$. Его масса: $dm = \rho \pi r^2 dz = \rho \pi (R^2 - z^2) dz$. Момент инерции диска относительно его оси: $dI = \frac{1}{2} dm \cdot r^2 = \frac{1}{2}\rho\pi(R^2-z^2)^2 dz$.
\[
	I = \int_{-R}^{R} \frac{1}{2}\rho\pi(R^2 - z^2)^2 dz = \rho\pi \int_0^R (R^2 - z^2)^2 dz.
\]
Раскрываем:
\[
	\int_0^R (R^4 - 2R^2z^2 + z^4)dz = R^5 - \frac{2R^5}{3} + \frac{R^5}{5} = R^5\left(1 - \frac{2}{3} + \frac{1}{5}\right) = \frac{8R^5}{15}.
\]
\[
	I = \rho\pi \cdot \frac{8R^5}{15} = \frac{3m}{4\pi R^3} \cdot \pi \cdot \frac{8R^5}{15} = \boxed{\frac{2}{5}mR^2}.
\]

\subsubsection{Теорема Штейнера}

\begin{theorem}[Гюйгенса--Штейнера]
	Момент инерции тела относительно произвольной оси равен сумме момента инерции относительно параллельной оси через центр масс и произведения массы на квадрат расстояния между осями:
	\[
		I = I_c + md^2.
	\]
\end{theorem}

\begin{proof}
	Пусть ось $z$ проходит через центр масс, а параллельная ось $z'$ находится на расстоянии $d$. Для элемента $dm$ на расстоянии $r$ от оси $z$ и $r'$ от оси $z'$:
	\[
		r'^2 = (x - a)^2 + (y - b)^2 = x^2 + y^2 - 2ax - 2by + a^2 + b^2,
	\]
	где $(a,b)$~--- координаты оси $z'$ относительно $z$, $d^2 = a^2 + b^2$.
	\[
		I' = \int r'^2 dm = \int(x^2+y^2)dm - 2a\int x\,dm - 2b\int y\,dm + d^2\int dm.
	\]
	Т.\,к. ось $z$ через центр масс: $\int x\,dm = 0$, $\int y\,dm = 0$.
	\[
		I' = I_c + md^2.
	\]
\end{proof}

\subsubsection{Энергия вращательного движения}

\[
	E_{\text{вр}} = \sum_i \frac{m_i V_i^2}{2} = \sum_i \frac{m_i (\omega r_i)^2}{2} = \frac{\omega^2}{2} \sum_i m_i r_i^2 = \frac{I\omega^2}{2}.
\]

Для тела, которое одновременно катится и движется поступательно:
\[
	E = \frac{mV_c^2}{2} + \frac{I_c \omega^2}{2}.
\]

\subsubsection{Закон сохранения момента импульса}

\begin{definition}
	Момент импульса: $\vec{L} = I\vec{\omega}$ (для вращения вокруг неподвижной оси); для точечного тела: $L = mvr\sin\alpha = pr_\perp$.\\
	$[L] = \text{кг}\cdot\text{м}^2/\text{с}$.
\end{definition}

Из основного уравнения:
\[
	M = I\beta = I\frac{d\omega}{dt} = \frac{dL}{dt} \quad \Longrightarrow \quad M\Delta t = \Delta L.
\]

\begin{law}[сохранения момента импульса]
	Если суммарный момент внешних сил равен нулю, то момент импульса системы сохраняется:
	\[
		M = 0 \quad \Longrightarrow \quad L = I\omega = \mathrm{const}.
	\]
\end{law}

\newpage
%===========================================
%           МКТ и ТЕРМОДИНАМИКА
%===========================================
\section{Молекулярно-кинетическая теория и термодинамика}

\subsection{Основные положения МКТ}

\begin{definition}[Молекулярно-кинетическая теория]
	Раздел физики, объясняющий свойства макроскопических тел и тепловых процессов на основе представления о том, что тела состоят из отдельных беспорядочно движущихся частиц.
\end{definition}

\textbf{Три основных положения МКТ:}
\begin{enumerate}
	\item Все вещества состоят из молекул (атомов).
	\item Молекулы находятся в непрерывном хаотическом (тепловом) движении.
	\item Между молекулами действуют силы притяжения и отталкивания.
\end{enumerate}

\paragraph{Опытные обоснования.}
\begin{itemize}
	\item \textbf{Положение 1:} делимость вещества, электронная и ионная микроскопия, рассеяние рентгеновских лучей.
	\item \textbf{Положение 2:} диффузия, броуновское движение (опыт Броуна, 1827; теория Эйнштейна, 1905).
	\item \textbf{Положение 3:} существование жидкого и твёрдого состояний, поверхностное натяжение, упругость.
\end{itemize}

\paragraph{Массы и размеры молекул.}
Относительная молекулярная масса: $M_r = \dfrac{m_0}{\frac{1}{12}m_{0_C}}$.\\
Число Авогадро: $N_A = 6{,}022 \cdot 10^{23} \; \text{моль}^{-1}$.\\
Молярная масса: $M$~--- масса одного моля вещества, $m_0 = M/N_A$.\\
Количество вещества: $\nu = \dfrac{N}{N_A} = \dfrac{m}{M}$.

\paragraph{Диффузия}~--- самопроизвольное проникновение молекул одного вещества в другое, обусловленное тепловым движением.

\paragraph{Броуновское движение}~--- хаотическое движение малых частиц, взвешенных в жидкости или газе, вызванное ударами молекул среды. Является прямым доказательством теплового движения молекул.

\paragraph{Взаимодействие молекул.}
На больших расстояниях~--- притяжение ($\sim r^{-7}$), на малых~--- отталкивание. На расстоянии $r_0$ (равновесное)~--- сила равна нулю, потенциальная энергия минимальна.

\subsection{Свойства газов. Идеальный газ}

\begin{definition}[Модель идеального газа]
	\begin{itemize}
		\item Суммарный объём молекул пренебрежимо мал по сравнению с объёмом сосуда.
		\item Столкновения молекул со стенками абсолютно упругие.
		\item Время столкновения много меньше времени свободного пробега.
		\item Нет дальнодействующего взаимодействия между молекулами.
	\end{itemize}
\end{definition}

\subsubsection{Вывод основного уравнения МКТ}

Рассмотрим кубический сосуд со стороной $L$. Молекула массы $m_0$ с проекцией скорости $v_x$ ударяется о стенку и упруго отражается. Изменение импульса при одном ударе: $\Delta p = 2m_0 v_x$.

Время между последовательными ударами о ту же стенку: $\Delta t = \dfrac{2L}{v_x}$.

Средняя сила от одной молекулы: $f = \dfrac{\Delta p}{\Delta t} = \dfrac{m_0 v_x^2}{L}$.

Суммарная сила от $N$ молекул: $F = \sum_{i=1}^{N} \dfrac{m_0 v_{xi}^2}{L} = \dfrac{N m_0 \overline{v_x^2}}{L}$.

Из изотропии: $\overline{v_x^2} = \dfrac{1}{3}\overline{v^2}$.

Давление: $p = \dfrac{F}{L^2} = \dfrac{N m_0 \overline{v^2}}{3L^3} = \dfrac{N m_0 \overline{v^2}}{3V}$.

\[
	\boxed{p = \frac{1}{3} m_0 n \overline{v^2} = \frac{2}{3} n \overline{E_{\text{к}}}}, \qquad \text{где } n = N/V, \quad \overline{E_{\text{к}}} = \frac{m_0 \overline{v^2}}{2}.
\]

\paragraph{Средняя квадратичная скорость.}
$v_{\text{кв}} = \sqrt{\overline{v^2}} = \sqrt{\dfrac{3kT}{m_0}} = \sqrt{\dfrac{3RT}{M}}$.

\subsubsection{Температура}

\begin{definition}
	Температура~--- мера средней кинетической энергии хаотического (поступательного) движения молекул:
	\[
		\overline{E_{\text{к}}} = \frac{3}{2} kT, \qquad k = 1{,}38 \cdot 10^{-23} \; \frac{\text{Дж}}{\text{К}} \; \text{(постоянная Больцмана)}.
	\]
\end{definition}

Абсолютный нуль ($T = 0 \; \text{К} = -273{,}15\;{^\circ}\text{C}$)~--- температура, при которой прекращается тепловое движение молекул (классически). Недостижим по третьему началу термодинамики.

\subsubsection{Уравнение состояния идеального газа}

Из $p = nkT = \dfrac{N}{V}kT = \dfrac{\nu N_A k T}{V}$, обозначая $R = N_A k = 8{,}314 \; \dfrac{\text{Дж}}{\text{моль}\cdot\text{К}}$:

\begin{definition}[Уравнение Менделеева--Клапейрона]
	\[
		pV = \nu RT.
	\]
\end{definition}

\begin{definition}[Уравнение Клапейрона (для фиксированной массы газа)]
	\[
		\frac{p_1 V_1}{T_1} = \frac{p_2 V_2}{T_2}.
	\]
\end{definition}

\subsubsection{Изопроцессы}

\begin{center}
\begin{tabular}{|l|c|c|c|}
	\hline
	\textbf{Процесс} & \textbf{Const} & \textbf{Закон} & \textbf{Графики}\\
	\hline
	Изотермический & $T$ & $p_1V_1 = p_2V_2$ (Бойль--Мариотт) & Гипербола в $(p,V)$\\
	\hline
	Изохорный & $V$ & $p_1/T_1 = p_2/T_2$ (Шарль) & Прямая через начало в $(p,T)$\\
	\hline
	Изобарный & $p$ & $V_1/T_1 = V_2/T_2$ (Гей-Люссак) & Прямая через начало в $(V,T)$\\
	\hline
\end{tabular}
\end{center}

\subsection{Распределение Максвелла}

\subsubsection{Пространство скоростей}

Состояние каждой молекулы задаётся вектором скорости $(v_x, v_y, v_z)$. Каждой молекуле соответствует точка в пространстве скоростей. Число молекул, скорости которых лежат в элементе $dv_x\,dv_y\,dv_z$:
\[
	dN = N f(v_x, v_y, v_z) \, dv_x \, dv_y \, dv_z.
\]

\subsubsection{Распределение по проекции скорости}

\[
	f(v_x) = \sqrt{\frac{m_0}{2\pi kT}} \exp\left(-\frac{m_0 v_x^2}{2kT}\right).
\]
Это гауссиана (колоколообразная кривая), симметричная относительно $v_x = 0$. Максимум при $v_x = 0$. Площадь под графиком равна 1.

\subsubsection{Распределение по модулю скорости (распределение Максвелла)}

Переходя к сферическим координатам в пространстве скоростей (элемент объёма $4\pi v^2 dv$):
\[
	\frac{dN}{N} = F(v) \, dv, \qquad F(v) = 4\pi v^2 \left(\frac{m_0}{2\pi kT}\right)^{3/2} \exp\left(-\frac{m_0 v^2}{2kT}\right).
\]

\paragraph{Физический смысл площадей под графиком.}
\begin{itemize}
	\item $\int_0^\infty F(v)\,dv = 1$ (нормировка).
	\item $\int_{v_1}^{v_2} F(v)\,dv$~--- доля молекул со скоростями в интервале $(v_1, v_2)$.
\end{itemize}

\paragraph{Зависимость от температуры и массы.}
\begin{itemize}
	\item При увеличении $T$: максимум графика $F(v)$ смещается вправо и становится ниже (кривая расплывается).
	\item При увеличении $m_0$: максимум смещается влево и становится выше (кривая сужается).
\end{itemize}

\paragraph{Характерные скорости.}

\begin{itemize}
	\item \textbf{Наиболее вероятная:} $v_{\text{нв}} = \sqrt{\dfrac{2kT}{m_0}} = \sqrt{\dfrac{2RT}{M}}$ (из $F'(v) = 0$).
	\item \textbf{Средняя:} $\langle v \rangle = \sqrt{\dfrac{8kT}{\pi m_0}} = \sqrt{\dfrac{8RT}{\pi M}} \approx 1{,}13 \, v_{\text{нв}}$.
	\item \textbf{Среднеквадратичная:} $v_{\text{кв}} = \sqrt{\overline{v^2}} = \sqrt{\dfrac{3kT}{m_0}} = \sqrt{\dfrac{3RT}{M}} \approx 1{,}22 \, v_{\text{нв}}$.
\end{itemize}

Соотношение: $v_{\text{нв}} < \langle v \rangle < v_{\text{кв}}$.

\paragraph{Опыт Штерна (1920).}
Серебро испарялось из нагретой нити внутри двух коаксиальных цилиндров. При вращении цилиндров быстрые и медленные атомы оседали на разных участках внешнего цилиндра, позволяя экспериментально определить распределение скоростей. Результаты подтвердили распределение Максвелла.

\subsection{Насыщенный и ненасыщенный пар. Влажность}

\begin{definition}[Насыщенный пар]
	Пар, находящийся в динамическом равновесии со своей жидкостью (скорость испарения = скорость конденсации). Его давление и плотность не зависят от объёма, но растут с температурой.
\end{definition}

\begin{definition}[Ненасыщенный пар]
	Пар, давление которого меньше давления насыщенного пара при данной температуре. Возможно дальнейшее испарение.
\end{definition}

\paragraph{Кипение}~--- интенсивное парообразование, происходящее по всему объёму жидкости (не только с поверхности). Условие кипения: давление насыщенного пара равно внешнему давлению. Температура кипения растёт с увеличением внешнего давления.

\paragraph{Критическая температура}
$T_{\text{кр}}$~--- температура, выше которой вещество не может существовать в жидком состоянии ни при каком давлении. При $T > T_{\text{кр}}$ различие между жидкостью и газом исчезает.

\paragraph{Критическое состояние}~--- точка на фазовой диаграмме ($T_{\text{кр}}, p_{\text{кр}}$), в которой жидкость и пар становятся тождественными по свойствам.

\paragraph{Влажность воздуха.}

\begin{definition}[Абсолютная влажность]
	$\rho_{\text{абс}} = \dfrac{m_{H_2O}}{V}$~--- плотность (масса) водяного пара в единице объёма воздуха.
\end{definition}

\begin{definition}[Относительная влажность]
	\[
		\varphi = \frac{\rho_{\text{абс}}}{\rho_{\text{нп}}(t)} \cdot 100\% = \frac{p_{H_2O}}{p_{\text{нп}}(t)} \cdot 100\%.
	\]
\end{definition}

\begin{definition}[Точка росы]
	Температура, при которой водяной пар в воздухе становится насыщенным (при данной абсолютной влажности): $\rho_{\text{абс}} = \rho_{\text{нп}}(t_{\text{росы}})$.
\end{definition}

\paragraph{Изотермы Эндрюса.}
Эндрюс (1869) экспериментально снял изотермы $CO_2$ в координатах $(V, p)$. При $T < T_{\text{кр}}$ изотерма имеет горизонтальный участок (двухфазная область: жидкость + пар). При $T = T_{\text{кр}}$~--- точка перегиба. При $T > T_{\text{кр}}$~--- монотонно убывающая кривая (газ не может быть сконденсирован).

\subsection{Реальные газы. Уравнение Ван-дер-Ваальса}

\begin{definition}[Уравнение Ван-дер-Ваальса (1873)]
	Учёт конечного объёма молекул и сил притяжения:
	\[
		\left(p + \frac{a\nu^2}{V^2}\right)(V - \nu b) = \nu RT,
	\]
	или для одного моля:
	\[
		\left(p + \frac{a}{V_m^2}\right)(V_m - b) = RT.
	\]
	Здесь $a$~--- параметр, учитывающий притяжение; $b$~--- параметр, учитывающий собственный объём молекул ($b = 4N_A V_0$, где $V_0$~--- объём одной молекулы).
\end{definition}

\paragraph{Смысл поправок.}
\begin{itemize}
	\item $\dfrac{a}{V_m^2}$~--- внутреннее давление: молекулы у стенки испытывают притяжение со стороны молекул внутри, что уменьшает давление на стенку.
	\item $b$~--- объём, недоступный для движения молекул из-за конечного размера самих молекул.
\end{itemize}

\paragraph{Семейство изотерм реального газа.}
При $T > T_{\text{кр}}$~--- кривая похожа на изотерму идеального газа. При $T < T_{\text{кр}}$~--- изотерма Ван-дер-Ваальса имеет S-образный участок (волнообразную кривую). Физически реализуется горизонтальный участок (правило Максвелла равных площадей), соответствующий фазовому переходу жидкость--пар.

\subsection{Фазовые переходы. Диаграмма состояния вещества}

\begin{definition}[Фаза]
	Термодинамически равновесное состояние вещества, отличающееся по физическим свойствам от других возможных равновесных состояний того же вещества.
\end{definition}

\begin{definition}[Фазовый переход]
	Превращение одной фазы в другую при изменении внешних условий (температуры, давления).
	\begin{itemize}
		\item \textbf{Фазовые переходы I рода:} сопровождаются выделением/поглощением теплоты и изменением объёма (плавление, кристаллизация, испарение, конденсация, сублимация, десублимация).
		\item \textbf{Фазовые переходы II рода:} без скрытой теплоты, но с изменением теплоёмкости, сжимаемости и т.\,п. (переход в сверхпроводящее состояние, переход ферромагнетик--парамагнетик).
	\end{itemize}
\end{definition}

\paragraph{Диаграмма состояния (p, T).}

На диаграмме три кривые:
\begin{itemize}
	\item Кривая плавления: граница твёрдое тело--жидкость.
	\item Кривая испарения (кипения): граница жидкость--пар, заканчивается в критической точке.
	\item Кривая сублимации: граница твёрдое тело--пар.
\end{itemize}

Все три кривые пересекаются в \textbf{тройной точке}~--- единственные $(p, T)$, при которых три фазы сосуществуют в равновесии. Для воды: $T = 273{,}16 \; \text{К}$, $p = 611 \; \text{Па}$.

Особенность воды: кривая плавления наклонена влево (при увеличении давления температура плавления уменьшается, т.\,к. лёд легче воды).

\subsection{Поверхностные явления}

\begin{definition}[Поверхностная энергия]
	$W_{\text{пов}} = \sigma \cdot S$, где $S$~--- площадь поверхности.
\end{definition}

\begin{definition}[Коэффициент поверхностного натяжения]
	$\sigma = \dfrac{W_{\text{пов}}}{S}$, $[\sigma] = \dfrac{\text{Дж}}{\text{м}^2} = \dfrac{\text{Н}}{\text{м}}$.\\
	Также: $\sigma = \dfrac{F}{l}$, где $l$~--- длина линии, вдоль которой действует сила поверхностного натяжения.
\end{definition}

\begin{definition}[Сила поверхностного натяжения]
	$F = \sigma l$, где $l$~--- длина контура (периметр смачивания).
\end{definition}

Поверхностное натяжение обусловлено тем, что молекулы на поверхности имеют меньше соседей, чем в объёме, поэтому обладают избыточной потенциальной энергией. Жидкость стремится минимизировать площадь поверхности (отсюда сферическая форма капель).

\paragraph{Смачивание и несмачивание.}
\begin{itemize}
	\item Смачивание~--- жидкость растекается по поверхности (краевой угол $\theta < 90^\circ$). Мениск вогнутый (вода в стекле).
	\item Несмачивание~--- жидкость собирается в каплю (краевой угол $\theta > 90^\circ$). Мениск выпуклый (ртуть в стекле).
\end{itemize}

\paragraph{Избыточное давление Лапласа.}
Под искривлённой поверхностью жидкости существует избыточное давление:
\[
	\Delta p = \frac{2\sigma}{R},
\]
где $R$~--- радиус кривизны поверхности. Для сферической капли~--- давление внутри больше, чем снаружи. Для пузыря (две поверхности): $\Delta p = \dfrac{4\sigma}{R}$.

\paragraph{Капиллярные явления.}
Высота подъёма жидкости в капилляре радиуса $r$:
\[
	h = \frac{2\sigma \cos\theta}{\rho g r},
\]
где $\theta$~--- краевой угол. При полном смачивании ($\theta = 0$): $h = \dfrac{2\sigma}{\rho g r}$.

\subsection{Тепловое расширение}

\begin{definition}[Линейное расширение]
	$l = l_0(1 + \alpha \Delta T)$, где $\alpha$~--- коэффициент линейного расширения, $[\alpha] = \text{К}^{-1}$.
\end{definition}

\begin{definition}[Объёмное расширение]
	$V = V_0(1 + \beta \Delta T)$, где $\beta \approx 3\alpha$~--- коэффициент объёмного расширения.
\end{definition}

Изменение плотности при нагревании:
\[
	\rho = \frac{\rho_0}{1 + \beta\Delta T} \approx \rho_0(1 - \beta\Delta T).
\]

\paragraph{Закон Гука через модуль Юнга.}
$\sigma = E \varepsilon$, т.\,е. $\dfrac{F}{S} = E \dfrac{\Delta l}{l_0}$, откуда $F = \dfrac{ES}{l_0}\Delta l = k\Delta l$, где $k = \dfrac{ES}{l_0}$.

\subsection{Внутренняя энергия. Первый закон термодинамики}

\begin{definition}[Внутренняя энергия]
	Суммарная кинетическая энергия хаотического движения молекул + потенциальная энергия их взаимодействия.

	Для идеального газа (нет потенциальной энергии взаимодействия):
	\[
		U = N \overline{E} = \frac{i}{2} \nu RT,
	\]
	где $i$~--- число степеней свободы молекулы.
\end{definition}

\begin{definition}[Степени свободы]
	\begin{itemize}
		\item Одноатомный газ: $i = 3$ (3 поступательных).
		\item Двухатомный газ: $i = 5$ (3 поступательных + 2 вращательных).
		\item Многоатомный ($\geqslant 3$ атомов): $i = 6$ (3 поступательных + 3 вращательных).
	\end{itemize}
\end{definition}

\textbf{Два способа изменения внутренней энергии:}
\begin{enumerate}
	\item Совершение работы над телом (или телом).
	\item Теплопередача: конвекция, теплопроводность, излучение.
\end{enumerate}

\begin{definition}[Работа газа]
	$A = p\Delta V$. В общем случае: $A = \int_{V_1}^{V_2} p \, dV$~--- площадь под графиком $p(V)$.
\end{definition}

\begin{law}[Первое начало термодинамики]
	Количество теплоты, переданное системе, идёт на изменение внутренней энергии и совершение работы:
	\[
		Q = \Delta U + A_{\text{газа}} \qquad \text{или} \qquad \Delta U = Q - A_{\text{газа}} = Q + A_{\text{над газом}}.
	\]
\end{law}

\paragraph{Применение первого начала к изопроцессам.}

\begin{center}
\begin{tabular}{|l|c|l|}
	\hline
	\textbf{Процесс} & \textbf{Первое начало} & \textbf{Комментарий}\\
	\hline
	Изохорный ($V = \text{const}$) & $Q = \Delta U = \dfrac{i}{2}\nu R\Delta T$ & $A = 0$\\
	\hline
	Изотермический ($T = \text{const}$) & $Q = A$ & $\Delta U = 0$\\
	\hline
	Изобарный ($p = \text{const}$) & $Q = \Delta U + p\Delta V = \dfrac{i+2}{2}\nu R\Delta T$ & $A = p\Delta V = \nu R\Delta T$\\
	\hline
	Адиабатный ($Q = 0$) & $0 = \Delta U + A$, т.\,е. $A = -\Delta U$ & При расширении $T$ убывает\\
	\hline
\end{tabular}
\end{center}

\subsubsection{Теплоёмкость}

\begin{definition}
	Удельная теплоёмкость: $c = \dfrac{Q}{m\Delta T}$, $[c] = \dfrac{\text{Дж}}{\text{кг}\cdot\text{К}}$.\\
	Молярная теплоёмкость: $C = c \cdot M = \dfrac{Q}{\nu \Delta T}$.
\end{definition}

\begin{itemize}
	\item $C_V = \dfrac{i}{2}R$~--- молярная теплоёмкость при постоянном объёме.
	\item $C_p = \dfrac{i+2}{2}R = C_V + R$~--- молярная теплоёмкость при постоянном давлении.
\end{itemize}

\begin{definition}[Соотношение Майера]
	$C_p - C_V = R$.
\end{definition}

\begin{definition}[Показатель адиабаты (коэффициент Пуассона)]
	$\gamma = \dfrac{C_p}{C_V} = \dfrac{i+2}{i}$.\\
	Для одноатомного: $\gamma = 5/3$; двухатомного: $\gamma = 7/5$; многоатомного: $\gamma = 8/6 = 4/3$.
\end{definition}

\subsubsection{Адиабатный процесс}

\begin{definition}
	Адиабатный процесс~--- термодинамический процесс без теплообмена с окружающей средой ($Q = 0$). Реализуется в теплоизолированных системах или при очень быстрых процессах.
\end{definition}

\begin{statement}[Уравнение адиабаты (уравнение Пуассона)]
	\[
		pV^\gamma = \text{const}, \qquad TV^{\gamma-1} = \text{const}, \qquad T^\gamma p^{1-\gamma} = \text{const}.
	\]
\end{statement}

Работа газа при адиабатном расширении:
\[
	A = -\Delta U = \frac{i}{2}\nu R(T_1 - T_2) = \frac{p_1V_1 - p_2V_2}{\gamma - 1}.
\]

\subsection{Тепловые машины. Второе начало термодинамики}

\begin{definition}[Тепловая машина]
	Устройство, циклически преобразующее теплоту в работу. Состоит из: нагревателя (температура $T_1$), рабочего тела, холодильника (температура $T_2$).
\end{definition}

\begin{definition}[КПД тепловой машины]
	\[
		\eta = \frac{A}{Q_{\text{н}}} = \frac{Q_{\text{н}} - |Q_{\text{х}}|}{Q_{\text{н}}} = 1 - \frac{|Q_{\text{х}}|}{Q_{\text{н}}}.
	\]
\end{definition}

\begin{theorem}[Карно]
	Максимальный КПД тепловой машины, работающей между температурами $T_1$ и $T_2$ ($T_1 > T_2$):
	\[
		\eta_{\max} = \eta_{\text{Карно}} = 1 - \frac{T_2}{T_1} = \frac{T_1 - T_2}{T_1}.
	\]
	Этот КПД достигается только в обратимом цикле Карно (две изотермы + две адиабаты).
\end{theorem}

\paragraph{Холодильная машина.}
Работает в обратном цикле: забирает теплоту $Q_{\text{х}}$ от холодного тела и передаёт $Q_{\text{н}}$ горячему, затрачивая работу $A$.
\begin{definition}[Холодильный коэффициент]
	\[
		\xi = \frac{Q_{\text{х}}}{A} = \frac{T_2}{T_1 - T_2}.
	\]
\end{definition}

\subsubsection{Второе начало термодинамики}

\begin{statement}[Формулировка Клаузиуса]
	Невозможен процесс, единственным результатом которого является передача теплоты от холодного тела к горячему.
\end{statement}

\begin{statement}[Формулировка Кельвина (Томсона)]
	Невозможен циклический процесс, единственным результатом которого является полное превращение теплоты, полученной от одного источника, в работу (невозможен вечный двигатель второго рода).
\end{statement}

\paragraph{Доказательство эквивалентности формулировок.}

\textit{Кельвин $\Rightarrow$ Клаузиус} (от противного): Допустим, нарушается формулировка Клаузиуса, т.\,е. существует процесс, передающий теплоту $Q$ от холодного тела к горячему без каких-либо других изменений. Тогда используем обычную тепловую машину, работающую между этими телами: она берёт $Q_{\text{н}}$ у горячего, отдаёт $Q_{\text{х}}$ холодному и совершает работу $A = Q_{\text{н}} - Q_{\text{х}}$. Нарушающий процесс возвращает $Q_{\text{х}}$ от холодного к горячему. Итого: горячее тело отдало $Q_{\text{н}} - Q_{\text{х}}$, получено $A = Q_{\text{н}} - Q_{\text{х}}$ работы, холодное тело не изменилось~--- нарушена формулировка Кельвина. Противоречие.

\textit{Клаузиус $\Rightarrow$ Кельвин} (от противного): Допустим, существует машина, полностью превращающая теплоту $Q$ от горячего тела в работу $A = Q$. Используем эту работу в холодильнике, который передаёт теплоту от холодного тела к горячему. Итого: теплота перешла от холодного к горячему без других изменений~--- нарушена формулировка Клаузиуса. Противоречие.

\newpage
%===========================================
%           ЭЛЕКТРИЧЕСТВО
%===========================================
\section{Электричество}

\subsection{Электрический заряд. Закон Кулона}

\begin{definition}[Электрический заряд]
	Физическая величина, характеризующая свойство тел участвовать в электромагнитном взаимодействии. Единица: Кл. Заряд электрона: $e = 1{,}6 \cdot 10^{-19}$ Кл.
	
	Заряд квантован (кратен $e$), существует два знака заряда.
\end{definition}

\begin{definition}[Электризация тел]
	Процесс сообщения телу электрического заряда. Способы: трение, контакт, индукция.
\end{definition}

\begin{law}[сохранения электрического заряда]
	В замкнутой системе алгебраическая сумма зарядов остаётся постоянной:
	\[
		\sum_i q_i = \text{const}.
	\]
\end{law}

\begin{law}[Кулона (1785)]
	Сила взаимодействия двух точечных неподвижных зарядов:
	\[
		F = k\frac{|q_1||q_2|}{\varepsilon r^2}, \qquad k = \frac{1}{4\pi\varepsilon_0} = 9 \cdot 10^9 \; \frac{\text{Н}\cdot\text{м}^2}{\text{Кл}^2},
	\]
	$\varepsilon_0 = 8{,}85\cdot10^{-12} \; \text{Ф/м}$~--- электрическая постоянная, $\varepsilon$~--- диэлектрическая проницаемость среды.
	
	Условия: точечные заряды, неподвижные.
\end{law}

\begin{statement}[Принцип суперпозиции]
	Сила, действующая на заряд со стороны системы зарядов, равна векторной сумме сил, действующих со стороны каждого заряда в отдельности. Присутствие других зарядов не влияет на попарное взаимодействие.
\end{statement}

\subsection{Электрическое поле}

\begin{definition}[Электрическое поле]
	Особая форма материи, посредством которой осуществляется взаимодействие электрических зарядов. Идея Фарадея: заряды не действуют друг на друга непосредственно; каждый заряд создаёт вокруг себя поле, и это поле действует на другие заряды.
\end{definition}

\begin{definition}[Напряжённость электрического поля]
	\[
		\vec{E} = \frac{\vec{F}}{q_0},
	\]
	где $q_0$~--- пробный заряд (малый положительный). $[E] = \text{В/м} = \text{Н/Кл}$.
	
	Для точечного заряда $Q$:
	\[
		E = \frac{kQ}{\varepsilon r^2} = \frac{Q}{4\pi\varepsilon_0\varepsilon r^2}.
	\]
\end{definition}

\paragraph{Линии напряжённости.}
\begin{itemize}
	\item Начинаются на положительных зарядах, заканчиваются на отрицательных (или уходят в бесконечность).
	\item Не пересекаются.
	\item Густота линий пропорциональна модулю напряжённости.
\end{itemize}

\subsubsection{Теорема Гаусса}

\begin{definition}[Поток вектора напряжённости]
	$\Phi_E = \oint \vec{E} \cdot d\vec{S} = \oint E \cos\theta \, dS$.
\end{definition}

\begin{theorem}[Гаусса]
	Поток вектора напряжённости через замкнутую поверхность равен:
	\[
		\Phi_E = \oint \vec{E} \cdot d\vec{S} = \frac{\sum q_{\text{внутри}}}{\varepsilon_0} = 4\pi k \sum q_{\text{внутри}}.
	\]
\end{theorem}

\paragraph{Применения теоремы Гаусса.}

\subparagraph{Бесконечная равномерно заряженная плоскость} (поверхностная плотность заряда $\sigma$).

Гауссова поверхность~--- цилиндр, пересекающий плоскость перпендикулярно. Поток через два основания: $2ES = \dfrac{\sigma S}{\varepsilon_0}$.
\[
	\boxed{E = \frac{\sigma}{2\varepsilon_0}} = 2\pi k\sigma.
\]

\subparagraph{Поле между двумя бесконечными параллельными разноимённо заряженными пластинами} (плоский конденсатор, $+\sigma$ и $-\sigma$):
\[
	E = \frac{\sigma}{\varepsilon_0} \quad \text{(между пластинами)}, \qquad E = 0 \quad \text{(вне пластин)}.
\]

\paragraph{Теорема Ирншоу.}
Система неподвижных точечных зарядов не может находиться в устойчивом равновесии только под действием электростатических сил. (Следствие теоремы Гаусса: в области без зарядов потенциал не имеет экстремумов, а значит, нет устойчивого положения равновесия.)

\subsection{Работа сил электрического поля. Потенциал}

Электростатическое поле~--- потенциальное (консервативное): работа сил поля не зависит от пути, а зависит только от начального и конечного положений заряда.

\begin{definition}[Работа в однородном поле]
	$A = qEl\cos\alpha$, где $l$~--- перемещение, $\alpha$~--- угол между $\vec{E}$ и $\vec{l}$.
\end{definition}

\begin{definition}[Работа в поле точечного заряда]
	\[
		A = kQq\left(\frac{1}{r_1} - \frac{1}{r_2}\right).
	\]
\end{definition}

\begin{definition}[Потенциальная энергия заряда]
	В поле точечного заряда: $W = \dfrac{kQq}{r}$.\\
	В однородном поле: $W = qEd$ (где $d$~--- расстояние вдоль поля).
\end{definition}

\begin{definition}[Потенциал]
	$\varphi = \dfrac{W}{q}$, $[\varphi] = \text{В}$ (Вольт).\\
	Потенциал точечного заряда: $\varphi = \dfrac{kQ}{r}$.\\
	Принцип суперпозиции: $\varphi = \sum_i \varphi_i$.
\end{definition}

\begin{definition}[Разность потенциалов (напряжение)]
	$U = \varphi_1 - \varphi_2 = \dfrac{A}{q}$.
\end{definition}

\begin{definition}[Эквипотенциальные поверхности]
	Поверхности, на которых $\varphi = \text{const}$. Перпендикулярны линиям напряжённости. Работа при перемещении заряда по эквипотенциальной поверхности равна нулю.
\end{definition}

\begin{statement}[Связь напряжённости и потенциала]
	В однородном поле: $E = \dfrac{U}{d} = \dfrac{\varphi_1 - \varphi_2}{d}$, где $d$~--- расстояние между эквипотенциалями вдоль силовой линии.\\
	В общем случае: $\vec{E} = -\nabla\varphi = -\text{grad}\,\varphi$, т.\,е. $E_x = -\dfrac{\partial\varphi}{\partial x}$ и т.\,д.
\end{statement}

\subsection{Электрическая ёмкость. Конденсаторы}

\begin{definition}[Электроёмкость]
	$C = \dfrac{q}{U}$, $[C] = \text{Ф}$ (Фарад).
\end{definition}

\begin{definition}[Ёмкость плоского конденсатора]
	\[
		C = \frac{\varepsilon\varepsilon_0 S}{d},
	\]
	где $S$~--- площадь пластин, $d$~--- расстояние между ними, $\varepsilon$~--- диэлектрическая проницаемость.
\end{definition}

\paragraph{Соединения конденсаторов.}

\subparagraph{Параллельное:} $C_{\text{об}} = C_1 + C_2 + \ldots$; \quad $U_{\text{об}} = U_1 = U_2$; \quad $q_{\text{об}} = q_1 + q_2$.

\subparagraph{Последовательное:} $\dfrac{1}{C_{\text{об}}} = \dfrac{1}{C_1} + \dfrac{1}{C_2} + \ldots$; \quad $q_{\text{об}} = q_1 = q_2$; \quad $U_{\text{об}} = U_1 + U_2$.

\begin{definition}[Энергия заряженного конденсатора]
	\[
		W = \frac{CU^2}{2} = \frac{q^2}{2C} = \frac{qU}{2}.
	\]
\end{definition}

\begin{definition}[Объёмная плотность энергии электрического поля]
	\[
		w = \frac{\varepsilon\varepsilon_0 E^2}{2}.
	\]
\end{definition}

\subsection{Законы постоянного тока}

\begin{definition}[Электрический ток]
	Упорядоченное (направленное) движение заряженных частиц.
\end{definition}

Условия существования постоянного тока:
\begin{itemize}
	\item наличие свободных заряженных частиц;
	\item наличие электрического поля (разности потенциалов);
	\item замкнутая цепь.
\end{itemize}

\begin{definition}[Сила тока]
	$I = \dfrac{q}{t} = \dfrac{dq}{dt}$, $[I] = \text{А}$.
\end{definition}

\begin{definition}[Плотность тока]
	$j = \dfrac{I}{S} = nev_d$, $[j] = \text{А/м}^2$,
	где $n$~--- концентрация носителей, $v_d$~--- дрейфовая скорость.
\end{definition}

\begin{definition}[Электродвижущая сила (ЭДС)]
	$\mathcal{E} = \dfrac{A_{\text{ст}}}{q}$, $[\mathcal{E}] = \text{В}$~--- работа сторонних сил по переносу единичного заряда по замкнутой цепи.
\end{definition}

\subsubsection{Законы Ома}

\begin{law}[Ома для участка цепи]
	$I = \dfrac{U}{R}$, где $R = \dfrac{\rho l}{S}$~--- сопротивление проводника.
\end{law}

\begin{law}[Ома для участка цепи с ЭДС]
	$I = \dfrac{\varphi_1 - \varphi_2 + \mathcal{E}}{R + r}$,
	где $r$~--- внутреннее сопротивление источника, $\mathcal{E}$ берётся со знаком <<+>>, если ЭДС помогает току.
\end{law}

\begin{law}[Ома для полной цепи]
	$I = \dfrac{\mathcal{E}}{R + r}$.
\end{law}

\paragraph{Короткое замыкание.} $R = 0$: $I_{\text{кз}} = \dfrac{\mathcal{E}}{r}$~--- максимальный ток, ограниченный только внутренним сопротивлением.

\paragraph{Соединения проводников.}

\subparagraph{Последовательное:}
$I = I_1 = I_2$, \quad $U = U_1 + U_2$, \quad $R = R_1 + R_2$.

\subparagraph{Параллельное:}
$I = I_1 + I_2$, \quad $U = U_1 = U_2$, \quad $\dfrac{1}{R} = \dfrac{1}{R_1} + \dfrac{1}{R_2}$.

\paragraph{Шунт к амперметру.}
Для расширения предела измерения в $n$ раз нужен шунт:
\[
	R_{\text{ш}} = \frac{R_A}{n - 1}.
\]

\paragraph{Добавочное сопротивление к вольтметру.}
Для расширения предела измерения в $n$ раз:
\[
	R_{\text{доб}} = (n-1)R_V.
\]

\subsubsection{Работа и мощность тока. Закон Джоуля--Ленца}

\begin{definition}[Работа тока]
	$A = IUt = I^2 Rt = \dfrac{U^2}{R}t$.
\end{definition}

\begin{definition}[Мощность тока]
	$P = IU = I^2R = \dfrac{U^2}{R}$.
\end{definition}

\begin{law}[Джоуля--Ленца]
	Количество теплоты, выделяемой в проводнике:
	$Q = I^2 R t$.
\end{law}

\subsubsection{Законы Кирхгофа}

\begin{law}[Первый закон Кирхгофа (правило узлов)]
	Алгебраическая сумма токов в узле равна нулю:
	\[
		\sum_i I_i = 0.
	\]
	(Следствие закона сохранения заряда.)
\end{law}

\begin{law}[Второй закон Кирхгофа (правило контуров)]
	В любом замкнутом контуре алгебраическая сумма ЭДС равна алгебраической сумме падений напряжения:
	\[
		\sum_i \mathcal{E}_i = \sum_i I_i R_i.
	\]
	(Следствие закона сохранения энергии и потенциальности электростатического поля.)
\end{law}

\subsection{Электрический ток в различных средах}

\subsubsection{Ток в металлах. Теория Друде}

\paragraph{Природа тока в металлах.}
Носители тока~--- свободные электроны (электронный газ), $n \sim 10^{28} \; \text{м}^{-3}$.

\paragraph{Скорость дрейфа и скорость распространения тока.}
Дрейфовая скорость электронов очень мала ($\sim 0{,}1$ мм/с), но электрический сигнал (электрическое поле) распространяется со скоростью, близкой к скорости света ($\sim 3 \cdot 10^8$ м/с), поскольку поле включает упорядоченное движение всех электронов почти одновременно.

\paragraph{Модель Друде.}
Электрон ускоряется полем $E$ между столкновениями. Среднее время свободного пробега: $\tau = \lambda / u$, где $\lambda$~--- средняя длина свободного пробега, $u = \sqrt{3kT/m_e}$~--- тепловая скорость.

Дрейфовая скорость: $v_d = \dfrac{eE\tau}{2m_e}$.

Плотность тока: $j = nev_d = \dfrac{ne^2\lambda}{2m_e u} E = \sigma_{\text{пр}} E$ (закон Ома в дифференциальной форме).

Удельная проводимость: $\sigma_{\text{пр}} = \dfrac{ne^2\lambda}{2m_e u}$; удельное сопротивление: $\rho = 1/\sigma_{\text{пр}}$.

\paragraph{Зависимость сопротивления металлов от температуры.}
$R = R_0(1 + \alpha T)$, где $\alpha > 0$~--- температурный коэффициент сопротивления. При нагревании колебания ионов решётки усиливаются, $\lambda$ уменьшается, сопротивление растёт.

\paragraph{Сверхпроводимость.}
При охлаждении некоторых материалов ниже критической температуры $T_c$ их сопротивление скачком обращается в нуль. Открыта Камерлинг-Оннесом (1911) для ртути при $T_c = 4{,}2$ К.

\subsubsection{Ток в растворах и расплавах электролитов}

\begin{definition}[Электролиты]
	Вещества, растворы или расплавы которых проводят ток благодаря наличию ионов (ионная проводимость).
\end{definition}

\begin{definition}[Электролитическая диссоциация]
	Распад молекул электролита на положительные и отрицательные ионы при растворении или плавлении.
\end{definition}

\begin{definition}[Электролиз]
	Окислительно-восстановительные процессы на электродах при прохождении тока через электролит.
\end{definition}

\begin{law}[Фарадея (электролиз)]
	Масса вещества, выделившегося на электроде:
	\[
		m = kIt = \frac{M}{zN_A e} It = \frac{M}{zF} It,
	\]
	где $k = \dfrac{M}{zF}$~--- электрохимический эквивалент, $z$~--- валентность, $F = eN_A = 96\,485$ Кл/моль~--- постоянная Фарадея.
\end{law}

\subsubsection{Ток в газах}

В обычных условиях газы~--- диэлектрики. Для протекания тока необходима ионизация.

\paragraph{Виды ионизации.}
\begin{itemize}
	\item \textbf{Термическая}~--- при высоких температурах.
	\item \textbf{Фотоионизация}~--- под действием света (УФ, рентген).
	\item \textbf{Ионизация электронным ударом}~--- ускоренные электроны выбивают электроны из атомов газа.
\end{itemize}

\begin{definition}[Несамостоятельный разряд]
	Ток в газе, существующий только при наличии внешнего ионизатора. При его удалении ток прекращается.
\end{definition}

\begin{definition}[Самостоятельный разряд]
	Ток в газе, поддерживаемый без внешнего ионизатора за счёт ионизации электронным ударом. Виды: искровой, дуговой, тлеющий, коронный.
\end{definition}

\paragraph{Плазма}~--- сильно ионизированный газ, в котором концентрации положительных и отрицательных зарядов практически равны (квазинейтральность).

\subsubsection{Ток в вакууме}

В вакууме нет носителей тока. Их создают \textbf{термоэлектронной эмиссией}~--- нагревом катода.

\begin{definition}[Вакуумный диод]
	Прибор с нагреваемым катодом и анодом в вакууме. Пропускает ток только в одном направлении (от катода к аноду при положительном напряжении на аноде). Используется для выпрямления тока.
\end{definition}

\subsubsection{Ток в полупроводниках}

\begin{definition}[Полупроводники]
	Вещества, электропроводность которых занимает промежуточное положение между проводниками и диэлектриками и сильно зависит от температуры, освещения, примесей. Примеры: Si, Ge, GaAs.
\end{definition}

\paragraph{Зависимость сопротивления от температуры.}
С ростом температуры сопротивление полупроводников \textit{уменьшается} (в отличие от металлов), т.\,к. растёт концентрация свободных носителей.

\paragraph{Собственная проводимость.}
При разрыве ковалентной связи возникает пара: свободный электрон + дырка. Оба типа носителей участвуют в проводимости.

\paragraph{Примесная проводимость.}

\begin{definition}[Доноры (n-тип)]
	Примесные атомы с большим числом валентных электронов (напр., P в Si). Дают лишние электроны. Основные носители~--- электроны.
\end{definition}

\begin{definition}[Акцепторы (p-тип)]
	Примесные атомы с меньшим числом валентных электронов (напр., B в Si). Создают дырки. Основные носители~--- дырки.
\end{definition}

\paragraph{p--n переход.}
На границе полупроводников p- и n-типа образуется область, обеднённая подвижными носителями (запирающий слой). Возникает контактная разность потенциалов. В прямом направлении (p к <<+>>, n к <<-->>)~--- ток проходит (запирающий слой сужается). В обратном~--- практически не проходит (запирающий слой расширяется).

\begin{definition}[Полупроводниковый диод]
	Прибор на основе p--n перехода, обладающий односторонней проводимостью.
\end{definition}

\begin{definition}[Транзистор]
	Прибор с двумя p--n переходами (p--n--p или n--p--n). Позволяет усиливать электрические сигналы. Малый ток базы управляет большим током коллектор--эмиттер.
\end{definition}

\newpage
%===========================================
%           МАГНЕТИЗМ
%===========================================
\section{Магнетизм}

\subsection{Магнитное поле}

\paragraph{Магнитное взаимодействие токов.}
Два параллельных проводника с токами в одном направлении притягиваются, в противоположных~--- отталкиваются. Это взаимодействие осуществляется посредством магнитного поля.

\begin{definition}[Магнитная индукция]
	Силовая характеристика магнитного поля. Определяется через максимальную силу Ампера:
	\[
		B = \frac{F_{\max}}{Il}, \qquad [B] = \text{Тл (Тесла)}.
	\]
\end{definition}

\paragraph{Линии магнитной индукции.}
\begin{itemize}
	\item Замкнуты (нет магнитных зарядов, $\mathrm{div}\,\vec{B} = 0$).
	\item Направление определяется правилом буравчика (правого винта).
\end{itemize}

\begin{definition}[Сила Ампера]
	Сила, действующая на проводник с током в магнитном поле:
	\[
		F_A = BIl\sin\alpha,
	\]
	где $\alpha$~--- угол между $\vec{B}$ и направлением тока. Направление~--- по правилу левой руки.
\end{definition}

\begin{definition}[Сила Лоренца]
	Сила, действующая на движущийся заряд в магнитном поле:
	\[
		F_L = qvB\sin\alpha,
	\]
	где $\alpha$~--- угол между $\vec{v}$ и $\vec{B}$.
\end{definition}

\paragraph{Движение заряженных частиц в однородном магнитном поле.}

Сила Лоренца перпендикулярна скорости, поэтому не совершает работы, а лишь изменяет направление скорости.

\begin{itemize}
	\item Если $\vec{v} \perp \vec{B}$: частица движется по окружности. Радиус:
	\[
		R = \frac{mv}{qB}, \qquad \text{период:} \quad T = \frac{2\pi m}{qB} \quad \text{(не зависит от скорости!)}.
	\]
	\item Если $\vec{v}$ составляет угол с $\vec{B}$: движение по винтовой линии (спирали).
\end{itemize}

\subsection{Закон Био--Савара--Лапласа}

\begin{law}[Био--Савара--Лапласа]
	Магнитная индукция $d\vec{B}$, создаваемая элементом тока $Id\vec{l}$ в точке на расстоянии $r$:
	\[
		dB = \frac{\mu_0 \mu}{4\pi} \frac{I \, dl \sin\alpha}{r^2},
	\]
	где $\alpha$~--- угол между $d\vec{l}$ и $\vec{r}$, $\mu_0 = 4\pi \cdot 10^{-7} \; \text{Гн/м}$~--- магнитная постоянная, $\mu$~--- магнитная проницаемость среды.
\end{law}

\paragraph{Поле бесконечного прямого проводника с током.}

Интегрируя закон Био--Савара--Лапласа по всей длине провода:
\[
	B = \frac{\mu_0 \mu I}{2\pi R_0},
\]
где $R_0$~--- расстояние от провода до точки наблюдения. Линии поля~--- концентрические окружности.

\paragraph{Поле конечного прямого проводника.}
Если проводник виден из точки наблюдения под углами $\alpha_1$ и $\alpha_2$ (от перпендикуляра):
\[
	B = \frac{\mu_0\mu I}{4\pi R_0}(\sin\alpha_1 + \sin\alpha_2).
\]
При $\alpha_1 = \alpha_2 = 90^\circ$ получаем формулу для бесконечного проводника.

\paragraph{Поле кругового витка с током на его оси.}
На расстоянии $x$ от центра витка радиуса $R$ вдоль оси, перпендикулярной плоскости витка:
\[
	B = \frac{\mu_0\mu I R^2}{2(R^2 + x^2)^{3/2}}.
\]
В центре витка ($x = 0$):
\[
	B = \frac{\mu_0\mu I}{2R}.
\]

\paragraph{Поле внутри бесконечного соленоида (катушки).}
Внутри однородное поле:
\[
	B = \mu_0 \mu n I = \frac{\mu_0\mu NI}{l},
\]
где $n = N/l$~--- число витков на единицу длины, $N$~--- общее число витков, $l$~--- длина соленоида.

\subsection{Электромагнитная индукция}

\begin{definition}[Магнитный поток]
	$\Phi = \vec{B} \cdot \vec{S} = BS\cos\theta$, где $\theta$~--- угол между $\vec{B}$ и нормалью $\vec{n}$ к поверхности.\\
	$[\Phi] = \text{Вб}$ (Вебер).
\end{definition}

\begin{law}[Фарадея (электромагнитная индукция)]
	ЭДС индукции равна скорости изменения магнитного потока:
	\[
		\mathcal{E}_{\text{инд}} = -\frac{d\Phi}{dt}.
	\]
	Для катушки из $N$ витков: $\mathcal{E}_{\text{инд}} = -N\dfrac{d\Phi}{dt}$.
\end{law}

\paragraph{Индукционное электрическое поле.}
Изменяющееся магнитное поле порождает вихревое электрическое поле (линии замкнуты, в отличие от электростатического).

\begin{statement}[Правило Ленца]
	Индукционный ток направлен так, чтобы своим магнитным полем противодействовать изменению магнитного потока, вызвавшему его.
\end{statement}

\paragraph{Токи Фуко (вихревые токи).}
Индукционные токи в сплошных проводниках. Вызывают нагрев, используются в индукционных печах, электромагнитных тормозах. Для уменьшения потерь сердечники трансформаторов набирают из тонких изолированных пластин.

\subsubsection{Самоиндукция}

Изменение тока в контуре приводит к изменению собственного магнитного потока и, следовательно, к возникновению ЭДС самоиндукции.

\begin{definition}[Индуктивность]
	$L = \dfrac{\Phi}{I}$, $[L] = \text{Гн}$ (Генри).\\
	Для соленоида: $L = \mu_0\mu n^2 V = \dfrac{\mu_0\mu N^2 S}{l}$.
\end{definition}

\begin{definition}[ЭДС самоиндукции]
	$\mathcal{E}_{\text{с}} = -L\dfrac{dI}{dt}$.
\end{definition}

\subsubsection{Энергия магнитного поля}

\begin{definition}[Энергия катушки]
	$W = \dfrac{LI^2}{2}$.
\end{definition}

\begin{definition}[Объёмная плотность энергии магнитного поля]
	\[
		w = \frac{B^2}{2\mu_0\mu} = \frac{\mu_0\mu H^2}{2},
	\]
	где $H = B/(\mu_0\mu)$~--- напряжённость магнитного поля.
\end{definition}

\subsection{Магнитные свойства вещества}

\begin{definition}[Магнитная проницаемость]
	$\mu = \dfrac{B}{B_0}$, где $B$~--- поле в веществе, $B_0$~--- поле в вакууме.
\end{definition}

\paragraph{Классификация веществ.}

\begin{itemize}
	\item \textbf{Диамагнетики} ($\mu < 1$, $\mu \approx 1$). Намагничиваются против внешнего поля (ослабляют его). Примеры: вода, медь, висмут, инертные газы, органика.

	\item \textbf{Парамагнетики} ($\mu > 1$, $\mu \approx 1$). Намагничиваются вдоль внешнего поля (незначительно усиливают). Примеры: алюминий, платина, соли железа и никеля, кислород.

	\item \textbf{Ферромагнетики} ($\mu \gg 1$, до $10^4$--$10^5$). Сильно намагничиваются вдоль поля. Примеры: Fe, Co, Ni, их сплавы.
\end{itemize}

\paragraph{Природа ферромагнетизма. Спин и домены.}

\begin{definition}[Спин]
	Собственный момент импульса элементарной частицы (квантовое свойство, не связано с вращением в классическом смысле). Спин электрона создаёт собственный магнитный момент.
\end{definition}

\begin{definition}[Домены]
	Макроскопические области спонтанной намагниченности в ферромагнетике. Внутри домена все магнитные моменты атомов (спины) выстроены параллельно. В отсутствие внешнего поля направления намагниченности разных доменов хаотичны, поэтому суммарная намагниченность нулевая.
\end{definition}

При наложении внешнего поля домены, направленные вдоль поля, растут за счёт других, а затем происходит поворот намагниченности оставшихся доменов.

\paragraph{Намагниченность.}
$\vec{J} = (\mu - 1)\dfrac{\vec{B}_0}{\mu_0}$ или $\vec{B} = \mu_0(\vec{H} + \vec{J})$, где $\vec{J}$~--- вектор намагниченности (магнитный момент единицы объёма).

\paragraph{Температура Кюри.}
Температура, выше которой ферромагнетик теряет ферромагнитные свойства и становится парамагнетиком (тепловое движение разрушает доменную структуру). Для железа: $T_C \approx 1043$ К.

\paragraph{Петля гистерезиса.}
Зависимость $B(H)$ при циклическом перемагничивании ферромагнетика. Кривая не совпадает при намагничивании и размагничивании (петля). Характерные величины:
\begin{itemize}
	\item \textbf{Остаточная индукция} $B_r$~--- значение $B$ при $H = 0$ (остаточная намагниченность).
	\item \textbf{Коэрцитивная сила} $H_c$~--- значение $H$, необходимое для полного размагничивания ($B = 0$).
\end{itemize}

Площадь петли гистерезиса пропорциональна потерям энергии за один цикл перемагничивания.

\begin{itemize}
	\item \textbf{Магнитомягкие материалы:} узкая петля (малые $B_r$, $H_c$), используются в сердечниках трансформаторов.
	\item \textbf{Магнитожёсткие материалы:} широкая петля (большие $B_r$, $H_c$), используются для постоянных магнитов.
\end{itemize}

\end{document}